\ExerciseSection

\begin{enumerate}
\def\labelenumi{\arabic{enumi})}
\item
  A linearly ordered group G (not necessarily commutative) is said to be Archimedean if the set I of elements \(\geqslant e\) (the identity element of G) satisfies axiom \((\mathrm{GR}_{\mathrm{IV}})\) of §2. Show that a linearly ordered group G is isomorphic to a subgroup of the additive group R if and only if G is Archimedean (distinguish two cases according as the set of elements \textgreater e in G has or has not a smallest element, and use Proposition 1 of §2).
\item
  Let G be a linearly ordered group (not necessarily commutative), with more than one element. Then the topology \(\mathcal{C}_{0}(G)\) (Chapter I, § 2, Exercise 5) is compatible with the group structure of G. If G is connected in this topology, show that G is isomorphic to the additive group R (use Exercise 7 of Chapter IV, § 2, and Proposition 2 of § 2).
\item
  Let G be a (not necessarily commutative) linearly ordered group. If G, endowed with the topology \(\mathcal{C}_{0}(G)\) , is locally compact and not discrete, then G is locally isomorphic to R, and the identity component of G is an open subgroup isomorphic to R (use Exercise 6 of Chapter IV, § 2, and Proposition 2 of § 2).
\item
  Let G be a topological group which satisfies the following conditions: \((R_{I})\) G is connected.
\end{enumerate}

\((R_{\mathrm{II}})\) The complement \(G^{*}\) of the identity element \(e\) of \(G\) is not connected.

Then there exists a continuous bijective homomorphism of G onto R (in other words, G is algebraically isomorphic to R and its topology is finer than that of R). The proof is in several steps:

\begin{enumerate}
\def\labelenumi{\alph{enumi})}
\item
  Let \((\mathbf{U}_i)_{1 \leqslant i \leqslant n}\) be a partition of \(\mathbf{G}^*\) into open sets in \(\mathbf{G}^* (n \geqslant 2)\) . Show that each \(\mathbf{U}_i\) is open in \(\mathbf{G}\) , that \(e\) lies in each \(\overline{\mathbf{U}}_i\) , and that \(\mathbf{G}\) is Hausdorff. Deduce that the closures \(\overline{\mathbf{U}}_i = \mathbf{U}_i \cup \{e\}\) in \(\mathbf{G}\) are connected (Chapter I, § 11, Exercise 4).
\item
  Let A be a connected component of \(U_{i}\) . Show that for each index \(j \neq i\) , we have \(A^{-1}\overline{U}_{j} = A^{-1}\) (observe that \(A^{-1}\overline{U}_{j}\) is connected and contains \(A^{-1}\) , and that \(A^{-1}\overline{U}_{j} \subset G^{*}\) if \(j \neq i\) ).
\item
  Show that n must be equal to 2 {[}for each index i = 1, 2, \ldots, n take a component \(A_{i}\) of \(U_{i}\) ; if \(j \neq i\) we have, by b), \(A_{i}^{-1}A_{j} \subset A_{i}^{-1}\) , \(A_{j}^{-1}A_{i} \subset A_{j}^{-1} \subset U_{j}^{-1}\) , whence \(A_{i}^{-1} \subset U_{j}\) {]}. Deduce that \(G^{*}\) has exactly two components A and B, that \(B = A^{-1}\) and \(\overline{A} = \complement(A^{-1})\) .
\item
  The relation \(yx^{-1} \in \overline{\mathbf{A}}\) is an order relation, which makes \(\mathbf{G}\) into a linearly ordered group (show that \(\overline{\mathbf{A}}^2 = \overline{\mathbf{A}}\) and that \(x\overline{\mathbf{A}}x^{-1} = \overline{\mathbf{A}}\) for all \(x \in \mathbf{G}\) ).
\item
  The topology \(\mathfrak{C}_0(\mathbf{G})\) is coarser than the given topology \(\mathfrak{C}\) on \(\mathbf{G}\) (observe that \(\mathbf{A}\) is open with respect to \(\mathfrak{C}\) ).
\end{enumerate}

Complete the proof by using Exercise 2 above and Chapter 1, § 11, no. 2, Proposition 4. {[}Cf. Chapter VI, § 1, Exercise 12 b).{]}

\begin{enumerate}
\def\labelenumi{\alph{enumi})}
\setcounter{enumi}{5}
\tightlist
\item
  Show that if in addition we suppose G to be either locally compact or locally connected, then G is isomorphic to R.
\end{enumerate}

* 5) Give an example of a topology on \(\mathbf{R}\) , compatible with the group structure, for which \(\mathbf{R}\) is connected, locally compact and locally connected, and for which the complement of \(\{o\}\) in \(\mathbf{R}\) is connected (use the fact that \(\mathbf{R}\) and \(\mathbf{R}^n\) are algebraically isomorphic). *

\begin{enumerate}
\def\labelenumi{\arabic{enumi})}
\setcounter{enumi}{5}
\tightlist
\item
  Let G be a topological group which satisfies the following conditions: \((\mathrm{LR}_{\mathbf{I}})\) G is Hausdorff and locally connected.
\end{enumerate}

\((\mathrm{LR}_{\mathbf{II}})\) There exists a connected neighbourhood \(\mathbf{U}\) of the identity element \(e\) of \(\mathbf{G}\) such that the complement of \(e\) in \(\mathbf{U}\) is not connected. Show that, under these conditions, \(\mathbf{G}\) is locally isomorphic to \(\mathbf{R}\) .

{[}Prove first that if \(\mathbf{V}\) is any connected neighbourhood of \(e\) in \(\mathbf{G}\) contained in \(\mathbf{U}\) , then \(\mathbf{V} \cap \complement\{e\}\) is not connected, that \(e\) lies in the closure of every component of \(\mathbf{V} \cap \complement\{e\}\) , and that \(\mathbf{V} \cap \complement\{e\}\) has exactly two components, by reasoning as in Exercise 4. Then take the connected neighbourhood \(\mathbf{V}\) sufficiently small and define a linear ordering on \(\mathbf{V}\) such that \(\mathfrak{C}_0(\mathbf{V})\) is coarser than the topology induced on \(\mathbf{V}\) by the topology of \(\mathbf{G}\) , and such that Proposition 2 of § 2 applies to the set \(\mathbf{E}\) of elements \(\geqslant e\) of \(\mathbf{V}\) .{]}

\begin{enumerate}
\def\labelenumi{\arabic{enumi})}
\setcounter{enumi}{6}
\tightlist
\item
  Let G be a non-discrete locally compact group with a countable base of open sets, and let \(V_{0}\) be a compact neighbourhood of the identity element e of G which contains no subgroup of G other than \(\{e\}\) .
\end{enumerate}

\begin{enumerate}
\def\labelenumi{\alph{enumi})}
\tightlist
\item
  Let \((x_{n})\) be a sequence of points of \(\mathbf{V}_{0}\) which converges to \(e\) . For each \(n\) , there exists an integer \(p(n) > 0\) such that \(x_{n}^{k} \in \mathbf{V}_{0}\) whenever \(k \leqslant p(n)\) and \(x_{n}^{p(n)+1} \notin \mathbf{V}_{0}\) , and we have \(\lim_{n \to \infty} p(n) = +\infty\) . Show that there exists a subsequence \((x_{k_{n}})\) of \((x_{n})\) such that, for every rational number \(r\) with \(|r| \leqslant 1\) , the sequence \((x_{k_{n}}^{[rp(k_{n})]})\) has a limit \(f(r)\) in \(\mathbf{V}_{0}\) (use the diagonal technique); and that if \(r, r'\) are rational numbers such that \(|r| \leqslant 1\) , \(|r'| \leqslant 1\) and \(|r + r'| \leqslant 1\) we have
\end{enumerate}

\[
f (r + r ^ {\prime}) = f (r) f (r ^ {\prime}).
\]

\begin{enumerate}
\def\labelenumi{\alph{enumi})}
\setcounter{enumi}{1}
\item
  Show that \(f\) is continuous in a neighbourhood of \(o\) in \(\mathbf{Q}\) . {[}Argue by contradiction: if there were a sequence \((r_n)\) of rational numbers tending to \(o\) and such that \(f(r_n)\) tended to an element \(y \neq e\) in \(V_0\) , show that we should have \(y^m \in V_0\) for all integers \(m\) .{]}
\item
  Deduce from b) that there exists a non-trivial continuous homomorphism of \(\mathbf{R}\) into \(\mathbf{G}\) (use Proposition 6 of §1).
\item
  Let H be a closed normal subgroup of G other than G itself, and suppose that there is a neighbourhood of the identity element in G/H which contains no non-trivial subgroup. Show that there exists a continuous homomorphism \(f: R \rightarrow G\) such that the composition
\end{enumerate}

\[
\mathbf {R} \xrightarrow {f} \mathbf {G} \to \mathbf {G} / \mathbf {H}
\]

is non-trivial.

\begin{enumerate}
\def\labelenumi{\arabic{enumi})}
\setcounter{enumi}{7}
\tightlist
\item
  Let G be a topological group and let H be a closed subgroup of G, contained in the centre of G, and such that there is a continuous surjective homomorphism \(\varphi: \mathbf{R} \to \mathbf{G} / \mathbf{H}\) . Show that \(\mathbf{G}\) is abelian {[}if \(a_0\) is an element of \(\mathbf{G}\) which is not in \(\mathbf{H}\) , observe that there exists a sequence \((a_n)\) of elements of \(\mathbf{G}\) and a sequence \((c_n)\) of elements of \(\mathbf{H}\) such that \(a_n = c_n a_{n+1}\) and that the subgroup of \(\mathbf{G}\) generated by \(\mathbf{H}\) and the \(a_n\) is dense in \(\mathbf{G}\) {]}.
\end{enumerate}

